% =========================================================
%  sections/03-cpmk.tex
%  Bagian 3 — Capaian Pembelajaran Mata Kuliah (CPMK)
% =========================================================

\section{Capaian Pembelajaran Mata Kuliah (CPMK)}

CPMK adalah kemampuan atau kompetensi spesifik yang diharapkan dikuasai
mahasiswa setelah menyelesaikan mata kuliah \NamaMK. Setiap CPMK diturunkan
dari CPL yang dibebankan pada mata kuliah ini.

\begin{table}[htbp]
\centering
\caption{Rumusan CPMK dan keterkaitannya dengan CPL}
\label{tab:cpmk}
\begin{tabularx}{\textwidth}{@{}C{2.0cm}Y C{3.6cm} C{1.8cm}@{}}
\toprule
\headrow \textbf{Kode} & \textbf{Rumusan CPMK} & \textbf{CPL Terkait} & \textbf{Ranah Bloom} \\
\midrule
\CPMK{1} &
Merancang \istilah{game loop} dan logika fisika sederhana. &
\CPL{03} (utama), \CPL{04} (pendukung) &
C2--C4 \\
\addlinespace
\CPMK{2} &
Menerapkan algoritma kecerdasan buatan (AI) dasar untuk NPC. &
\CPL{03} (utama), \CPL{04} (pendukung) &
C3--C5 \\
\addlinespace
\CPMK{3} &
Menggunakan \istilah{game engine} (Unity) untuk membangun level. &
\CPL{04} (utama), \CPL{03} (pendukung) &
C3--C4 \\
\addlinespace
\CPMK{4} &
Mengintegrasikan aset audio dan visual ke dalam interaksi \istilah{game}. &
\CPL{04} (utama), \CPL{03} (pendukung) &
C3--C5 \\
\bottomrule
\end{tabularx}
\end{table}

\subsection{Matriks Hubungan CPL--CPMK}
\label{sec:matriks-cpl-cpmk}

\begin{table}[htbp]
\centering
\caption{Matriks hubungan CPL dan CPMK}
\label{tab:matriks-cpl-cpmk}
\begin{tabularx}{\textwidth}{@{}L{3.8cm}@{\extracolsep{\fill}}C{2.4cm} C{2.4cm} C{2.4cm} C{2.4cm}@{}}
\toprule
\headrow & \textbf{\CPMK{1}} & \textbf{\CPMK{2}} & \textbf{\CPMK{3}} & \textbf{\CPMK{4}} \\
\headrow & \footnotesize \istilah{Game loop} \& fisika & \footnotesize AI NPC & \footnotesize Engine \& level & \footnotesize Aset audio-visual \\
\midrule
\CPL{03} --- Algoritma dan bahasa          & \Utama & \Utama & \Pendukung & \Pendukung \\
\CPL{04} --- Desain, implementasi, evaluasi & \Pendukung & \Pendukung & \Utama & \Utama \\
\bottomrule
\end{tabularx}
\tabsource{Keterangan: \textbf{Utama} = CPMK merupakan perwujudan langsung CPL;
Pendukung = CPMK memperkuat pencapaian CPL.}
\end{table}

\subsection*{Penjelasan Hubungan}

\begin{itemize}
  \item \textbf{\CPL{03} $\rightarrow$ \CPMK{1}:} mahasiswa merancang algoritma
        siklus input--update--render, \istilah{timestep}, gerak, dan
        \istilah{collision}, lalu menerapkannya dengan Python/Pygame.
  \item \textbf{\CPL{03} $\rightarrow$ \CPMK{2}:} mahasiswa menganalisis dan
        mengimplementasikan algoritma AI (\istilah{chase}, FSM, A*) pada bahasa
        Python tanpa bergantung pada fitur \istilah{black-box}
        \istilah{game engine}.
  \item \textbf{\CPL{04} $\rightarrow$ \CPMK{3}:} mahasiswa mendesain dan
        mengimplementasikan solusi level/prototipe pada Unity, lalu
        mengevaluasi hasilnya.
  \item \textbf{\CPL{04} $\rightarrow$ \CPMK{4}:} mahasiswa mengintegrasikan
        aset multimedia ke interaksi \istilah{game} dan mengevaluasi kualitas
        pengalaman audio-visual.
  \item \textbf{\CPL{04} juga mendukung \CPMK{1}--\CPMK{2}} pada tahap evaluasi
        prototipe sederhana (\istilah{playtest}, perbaikan \istilah{bug},
        refleksi desain).
  \item \textbf{\CPL{03} juga mendukung \CPMK{3}--\CPMK{4}} ketika mahasiswa
        memilih dan menyesuaikan algoritma/\istilah{script} C\# serta struktur
        data di dalam \istilah{engine}.
\end{itemize}

\subsection{Pendekatan Teknologi per CPMK}

\begin{table}[htbp]
\centering
\caption{Pendekatan teknologi yang digunakan pada setiap tahap}
\label{tab:teknologi}
\begin{tabularx}{\textwidth}{@{}L{4.6cm} L{5.0cm} Y@{}}
\toprule
\headrow \textbf{Tahap} & \textbf{Teknologi} & \textbf{CPMK} \\
\midrule
Fondasi algoritma   & Python 3 + Pygame & \CPMK{1}, \CPMK{2} \\
Implementasi \istilah{engine}     & Unity (C\#)        & \CPMK{3}, \CPMK{4} \\
\bottomrule
\end{tabularx}
\end{table}
